% Pursuing Stacks (cotes 134-2 to 134-8): archive cote 134-6, site page 1 (PDF page 2), typescript page 320.
% typescript. Draft compared: ps5.tex:1359-1431.
% AI-transcribed by Claude Opus 5.5 (claude-opus-5-5, adaptive thinking, effort high), Message Batches API, batch pilot (msgbatch_0189Nhf1xB615YWnztubtLfc), 2026-11-12; not yet checked by eye.
\scanpage{134-6/1}
\mspage{320}

5.7.

9\over{\unsure{1}}{2}. The review on ``pending questions and topics'' related to part IV
of the notes has taken pretty much longer than expected. It was quite
useful though, to get a clearer view of what those questions are about,
and to get a feeling for what to include and develop, and where. As I
do not intend to spend my life on the task, not even one year, it is
becoming clear that I am not going to get the whole picture of all the
questions touched - and some presumably I am going to leave just
aside, as they do not seem undispensable for a comprehensive overall
picture of what I\hand{'}m after. This seems to me to be the case for the
questions 7) and 8), concerned with finiteness conditions for homotopy
types in terms of models, and \xout{the question} with test categories with
boundary operations. At the opposite side, it seems that the questions
2), \inserted{part of} 3), and 9), about the notion of contractor, induced asphericity
and contractibility structures on a category $M_{/a}$ , and ``miscellaneous''
left-overs from part II, should be dealt with in part IV - whose end
now is in sight after all ! On the other hand, questions 1), another
part of 3), 5) and \xout{sixx} 6), about morphisms of asphericity structures
and related problems, about an autodual treatment of asphericity and
contractibility notions, about a \xout{caracterization of (certain)} \inserted[typed]{handy criterion for} canoni\hand{cal}
modelizers, and about existence theorems for various kinds of test
functors or aspheric functors, while I feel that I should come \xout{to} at
least to a considerably clearer understanding of these matters than now,
the adequate place for developping such reflection is definitely \uline{not} in
the present part IV, but belongs to \inserted[typed]{one or the other of} the three parts still ahead in our
overall reflection on the modelizing story.

During our review, we came a number of times upon situations where
the question arose as to whether \deleted{the} \inserted{one} point I like to make, namely that
a more or less arbitrary \hand{(}strict\hand{)} test category\xout{,} \xout{or contractor} ``is just
as good'' as the sacro-saint test category \hand{$\Delta$}, or it's twin brothers
\hand{$\square$} and \hand{\unsure{$\bigcirc$}}, is a valid one or not. I would like to list here these si-%
tuations, with a view of coming back to it later :

a) Existence theorems for test functors (cf section 90).

b) $A^{\wedge}$ and various other categories constructed in terms of A, such
as $\underline{\mathrm{Hom}}(A^{\circ},M )$, are closed model categories (under suitable as\xout{x}-%
sumptions ... ).

c) Independance of the derived category of $\underline{\mathrm{Hom}}(A^{\circ},M)$ on the choice
of test category, notably when M is a topos or the dual of a

\margin{1}
